\cvsection{Expériences professionnelles}

\begin{cventries}
  \cventry
    {Stagiaire ingénieur de recherche en IA}
    {INRIA}
    {Bordeaux, France}
    {Mars 2026 -- Septembre 2026}
    {
      \begin{cvitems}
        \item {Conception et développement en autonomie d’une application web fonctionnelle répondant au besoin de chercheurs souhaitant trouver des articles scientifiques.}
        \item {Développement d’un backend Python/FastAPI avec API REST documentées avec OpenAPI et persistance des documents et métadonnées dans MongoDB.}
        \item {Développement d’une interface React/Next.js connectée au backend pour permettre aux chercheurs d’interagir avec l’application, l’agent et ses outils.}
        \item {Mise en œuvre d’un pipeline de traitement de documents textuels avec chunking, embeddings, recherche sémantique et base vectorielle ChromaDB.}
        \item {Construction d’un agent multi-étapes sans framework d’orchestration avec gestion des outils, de l’état, de la mémoire et du contexte.}
        \item {Automatisation de la collecte de métadonnées d’articles scientifiques disponibles en ligne et transmission des résultats de tâches cron par webhooks.}
        \item {Gestion d’une application web conteneurisée permettant d’utiliser des LLM installés localement avec LM Studio sur un Mac Studio.}
        \item {Conteneurisation de composants avec Docker, versionnement avec Git et mise en place d’une CI/CD sur GitHub.}
      \end{cvitems}
    }

  \cventry
    {Stagiaire ingénieur de recherche en apprentissage par renforcement}
    {Université Karunya}
    {Coimbatore, Inde}
    {Juin 2025 -- Sept. 2025}
    {
      \begin{cvitems}
        \item {Développement d’un modèle d’apprentissage par renforcement et de son environnement d’entraînement pour la navigation autonome d’un robot hospitalier.}
        \item {Déploiement et optimisation du logiciel et du modèle sur Raspberry Pi pour des tests en conditions réelles.}
      \end{cvitems}
    }
\end{cventries}
